%% kapa: exam=1 round=37 year=2026 period=2 form=A subject=accounting subject_ko=회계학 seq=30 number=70
%% source: Q-Net 감정평가사 기출문제 (cst00309) artlSeq=5258924 file=raw/1008/5258924/2026년 제37회 감정평가사 1차 2교시.pdf
%% conversion: kordoc 4.24.51 -> tools/gen1.py (auto); review=proofread-sonnet; flags=none
\begin{question}{70}
(주)감평은 20×1년 초에 (주)한국의 발행주식 10 \%를 ₩50,000에 취득하여 당기손익-공정가치 측정 금융자산으로 분류하였다. 20×1년 말 이 금융자산의 공정가치는 ₩60,000이며, 20×2년 초 공정가치는 ₩62,000이다. 20×2년 초 (주)감평은 (주)한국의 나머지 주식 90 \%를 추가 취득하고 그 대가로 (주)감평 주식(주당 공정가치 ₩1,000, 주당액면금액 ₩500) 900주를 교부하여 흡수합병하였다. 합병일 현재 향후
(주)한국의 영업성과에 따라 추가지급하는 조건부대가의 공정가치는 ₩12,000이다. 합병당시 (주)한국의 재무상태표상 순자산(자산-부채) 장부금액은 ₩600,000이며, 공정가치는 ₩800,000이다. 합병일 현재 장부금액에 포함되어 있지 않은 개발프로젝트의 공정가치는 ₩20,000이며, 식별가능한 무형자산이다. 또한 합병일 현재 신뢰성 있게 측정한 우발부채의 공정가치는 ₩6,000이다. 합병일에 (주)감평이 인식할 영업권은? \par
\begin{choices}
  \item ₩160,000
  \item ₩154,000
  \item ₩152,000
  \item ₩140,000
  \item ₩120,000
\end{choices}
\end{question}
\answer{1}
